\section{Conclusion}
\label{sec:conclusion}

We opened this paper with a reversal: semantic entropy wins on TriviaQA by 0.155 AUROC, then loses on TruthfulQA by 0.066 AUROC.
That reversal is not a failure of the method --- it is a diagnostic signal revealing the condition under which semantic-level uncertainty estimation provides genuine value.

This paper provides the first unified comparison of SE, TE, SCG, and VC at Llama-2-7B scale, with direct mechanism measurement and cross-benchmark testing.
\textbf{We establish that SE substantially outperforms TE on TriviaQA} (AUROC 0.717 vs.\ 0.562, gap = +0.155), confirming the SE $>$ TE ordering is not scale-dependent noise.
\textbf{We identify the mechanism:} intra-cluster TE variance = 7.152 nats$^2$ (71$\times$ threshold) across 76/98 questions confirms TE aggregates paraphrase noise that SE's NLI clustering removes.
\textbf{We identify the task-structure scope condition:} the ordering reverses on TruthfulQA, where deterministic-wrong outputs defeat SE's filtering advantage.
\textbf{We characterize two alternative method failure modes:} BERTScore SCG diverges from SE by 0.336 AUROC on short QA, and VC at 7B scale produces degenerate confidence (ECE = 0.430).

\textbf{Future directions.}
(1) Re-run H-C1 with \texttt{nli-deberta-v3-large} to disentangle NLI model size from task structure.
(2) N=500 extension of H-E1 to confirm gap stability at larger N.
(3) Test SelfCheckNLI to determine whether NLI-based SCG recovers SE-equivalence on short QA.
(4) Scale to Llama-13B/70B to test whether the SE $>$ TE gap grows with scale.

The uncertainty method that wins on TriviaQA loses on TruthfulQA.
Rather than choosing between benchmarks, practitioners should ask what each task's incorrect outputs look like --- diverse or deterministic --- and select accordingly.
Understanding this scope condition is more valuable than an unconditional method ranking, and provides a foundation for principled uncertainty method selection as deployment settings grow more varied.
